\icode{E-MI-SI-L3-261029-002}
如图所示，竖直平面内有一圆环，圆心为\(O\)，半径为\(R\)，\(PQ\)为水平直径，\(MN\)为倾斜直径，\(PQ\)与\(MN\)间的夹角为\(\theta\)，一条不可伸长的轻绳长为\(L\)，两端分别固定在圆环的\(M\)、\(N\)两点，轻质滑轮连接\(M\)一个质量为\(m\)的重物，放置在轻绳上，不计滑轮与轻绳间的摩擦，滑轮大小不计，重力加速度为\(g\)。现将圆环从图示位置绕圆心\(O\)顺时针缓慢转过\(2\theta\)角，下列说法正确的是\paren[]
	
\begin{tikzpicture}[font=\scriptsize]
	\draw (0,0) circle(2) (-.5,0) arc(180:200:.5);
	\node at(-.8,-.15){\(\theta\)};
	\draw[dashed,cyan] (-2,0)node[left,black]{\(P\)}--(2,0)node[right,black]{\(Q\)};
    \draw[dashed,green] (200:2)node[left,black]{\(M\)}--(20:2)node[right,black]{\(N\)};
	\draw[dashed,teal] (-20:2)--(160:2);
	\draw[magenta] (200:2)--++(-30:{2*cos(10)/sin(60)-2*sin(10)/cos(60)})coordinate(A)--(20:2);
	\draw ($(A)+(0,{0.16*sqrt(3)/3})$)coordinate(O) circle(.08) (O)--++(0,-.3)coordinate(B);
	\draw[fill=orange!60] ($(B)+(-.08,0)$) rectangle($(B)+(.08,-.16)$);
	\fill (0,0)node[above]{\(O\)} circle(.05) (O) circle(.04);
\end{tikzpicture}
\begin{choices}
	\item 直径\(MN\)水平时，轻绳的张力大小为\(\frac{mgL}{\sqrt{L^2 - 4R^2}}\)
	\item 圆环从图示位置顺时针缓慢转过\(2\theta\)的过程中，轻绳的张角先减小再	增大
	\item 圆环从图示位置顺时针缓慢转过\(2\theta\)的过程中，重物的重力势能先增大后减小
	\item 圆环从图示位置顺时针缓慢转过\(\theta\)的过程中，轻绳的张力逐渐减小	
\end{choices}
\begin{answer}
	C\qquad 提示：动滑轮到\(PQ\)的距离为\(\sqrt{L^2 -(2R\cos\theta)^2}\)
\end{answer}